\documentclass[aspectratio=169, 10pt]{beamer}

% --- GÓI TIẾNG VIỆT & PHÔNG CHỮ ---
\usepackage[utf8]{inputenc}
\usepackage[vietnamese]{babel}
\usepackage{amsmath, amssymb}
\usepackage{booktabs}
\usepackage{tikz}
\usetikzlibrary{shapes.geometric, arrows, positioning}
\usepackage{colortbl}
\usepackage{xcolor}

% --- THIẾT KẾ GIAO DIỆN BEAMER (NAVY & CLEAN THEME) ---
\usetheme{Madrid}
\usecolortheme{whale}

\definecolor{primaryNavy}{RGB}{27, 54, 93}
\definecolor{secondarySlate}{RGB}{75, 107, 148}
\definecolor{lightBg}{RGB}{245, 247, 250}
\definecolor{accentGold}{RGB}{212, 160, 23}

\setbeamercolor{palette primary}{bg=primaryNavy, fg=white}
\setbeamercolor{palette secondary}{bg=secondarySlate, fg=white}
\setbeamercolor{palette tertiary}{bg=primaryNavy!90, fg=white}
\setbeamercolor{titlelike}{parent=palette primary}
\setbeamercolor{structure}{fg=primaryNavy}
\setbeamercolor{block title}{bg=primaryNavy, fg=white}
\setbeamercolor{block body}{bg=lightBg, fg=black}

\setbeamertemplate{navigation symbols}{}
\setbeamertemplate{footline}[frame number]

% --- THÔNG TIN BÁO CÁO ---
\title[Khảo sát hiện trạng - Hệ thống Đấu giá]{\textbf{PHÁT TRIỂN PHẦN MỀM HƯỚNG ĐỐI TƯỢNG}}
\subtitle{\textbf{Chương 3 (Phần 1): Khảo sát hiện trạng \& Phân loại yêu cầu}\\\small Đề tài: Hệ thống Sàn Đấu giá Trực tuyến Tự động}
\author[Nhóm 01]{Giảng viên phụ trách: \textbf{Phạm Thi Vương}\\Nhóm thực hiện: \textbf{Nhóm 01}}
\institute[OOAD]{Học phần: Phát Triển Phần Mềm Hướng Đối Tượng (OOAD)}
\date{\today}

\begin{document}

% ==================== SLIDE 1 ====================
\begin{frame}
    \titlepage
\end{frame}

% ==================== SLIDE 2 ====================
\begin{frame}{Nội dung Phần 1}
    \begin{columns}[t]
        \begin{column}{0.48\textwidth}
            \begin{block}{1. Giai đoạn khảo sát hiện trạng}
                Ba mảng cần tìm hiểu trước khi phân tích:\\
                Tổ chức -- Nghiệp vụ -- Tin học.
            \end{block}
            \vspace{0.3cm}
            \begin{block}{3. Hiện trạng Tin học}
                Hạ tầng máy chủ, công nghệ Java Spring Boot, MySQL và năng lực con người.
            \end{block}
        \end{column}
        
        \begin{column}{0.48\textwidth}
            \begin{block}{2. Ví dụ minh họa nghiệp vụ sàn}
                Quy trình Đấu giá ủy quyền tự động \& Chốt đơn, Ví bảo chứng thanh toán.
            \end{block}
            \vspace{0.3cm}
            \begin{block}{4. Xác định và phân loại yêu cầu}
                Chuyển từ khảo sát hiện trạng sang đặc tả yêu cầu (Chức năng \& Phi chức năng).
            \end{block}
        \end{column}
    \end{columns}
\end{frame}

% ==================== SLIDE 3 ====================
\begin{frame}
    \begin{center}
        \vspace{1.5cm}
        {\Large \textbf{\color{primaryNavy}PHẦN 1}}\\[0.5cm]
        {\huge \textbf{\color{primaryNavy}Giai đoạn khảo sát hiện trạng}}\\[0.6cm]
        \textit{\color{secondarySlate}Trước khi đặc tả yêu cầu, cần hiểu rõ hiện trạng tổ chức và nghiệp vụ của sàn}
    \end{center}
\end{frame}

% ==================== SLIDE 4 ====================
\begin{frame}{Ba mảng cần khảo sát}
    \begin{columns}[t]
        \begin{column}{0.32\textwidth}
            \begin{block}{\centering \textbf{Hiện trạng tổ chức}}
                Cơ cấu sàn giao dịch giữa người dùng với người dùng, tương tác giữa Ban quản trị, Kế toán và Khách hàng.
            \end{block}
        \end{column}
        \begin{column}{0.32\textwidth}
            \begin{block}{\centering \textbf{Hiện trạng nghiệp vụ}}
                Luồng tạo phiên, đấu giá ủy quyền tự động, giá sàn bảo lưu ẩn và thanh toán ví bảo chứng.
            \end{block}
        \end{column}
        \begin{column}{0.32\textwidth}
            \begin{block}{\centering \textbf{Hiện trạng Tin học}}
                Máy chủ, CSDL MySQL, giao thức đồng bộ thời gian thực và năng lực công nghệ người dùng.
            \end{block}
        \end{column}
    \end{columns}
\end{frame}

% ==================== SLIDE 5 ====================
\begin{frame}{Hiện trạng tổ chức (Sàn Đấu giá Trực tuyến)}
    \begin{itemize}
        \item \textbf{Đối nội -- Cơ cấu tổ chức nội bộ của sàn:}
        \begin{itemize}
            \item \textbf{Ban Quản trị sàn:} Giám sát toàn hệ thống, quản lý danh mục món hàng, áp dụng chế tài xử phạt người bùng hàng (\textit{Thẻ phạt bùng hàng}).
            \item \textbf{Bộ phận Kế toán / Tài chính:} Quản lý ví doanh thu sàn (\textit{Ví Doanh thu}), trích thu 5\% phí dịch vụ và kiểm soát ví ký quỹ trung gian (\textit{Ví Bảo chứng}).
            \item \textbf{Bộ phận Trọng tài khiếu nại:} Tiếp nhận bằng chứng khiếu nại, phân xử tranh chấp giữa Người mua và Người bán.
        \end{itemize}
        \vspace{0.3cm}
        \item \textbf{Đối ngoại -- Mối quan hệ với môi trường xung quanh:}
        \begin{itemize}
            \item \textbf{Người bán:} Đăng tải món hàng, thiết lập mức giá sàn bảo lưu ẩn.
            \item \textbf{Người mua:} Đặt mức giá trần tối đa bí mật, thanh toán và kiểm tra nhận hàng.
            \item \textbf{Đối tác thanh toán / Ngân hàng:} Thực hiện liên kết nạp/rút tiền điện tử an toàn.
        \end{itemize}
    \end{itemize}
\end{frame}

% ==================== SLIDE 6 ====================
\begin{frame}{Hiện trạng nghiệp vụ}
    \begin{itemize}
        \item \textbf{Mục tiêu số 1:} Thấu hiểu quy trình trả giá và tháo gỡ khó khăn của khách hàng (không phải túc trực canh màn hình liên tục, loại trừ rủi ro bùng hàng).
        \item \textbf{Các quy trình nghiệp vụ chính:}
        \begin{enumerate}
            \item Quản lý danh mục \& Đăng bán món hàng (\textit{Niêm yết sản phẩm}).
            \item Đấu giá ủy quyền tự động \& Mua đứt ngay.
            \item Chốt phiên tự động (\textit{Tiến trình hẹn giờ}) \& Khởi tạo Đơn hàng trúng giá.
            \item Thanh toán tạm giữ bảo chứng \& Tự động giải ngân.
            \item Xử lý vi phạm bùng hàng \& Hỗ trợ đăng lại phiên đấu giá.
        \end{enumerate}
        \item \textbf{Tần suất \& thời điểm thực hiện:}
        \begin{itemize}
            \item Trả giá: Diễn ra từng giây, đặc biệt dồn dập ở 5 phút cuối phiên.
            \item Quét kết thúc: Tự động mỗi 5 giây chạy ngầm trong hệ thống.
        \end{itemize}
    \end{itemize}
\end{frame}

% ==================== SLIDE 7 ====================
\begin{frame}{Hiện trạng nghiệp vụ -- Đánh giá \& Khó khăn}
    \begin{itemize}
        \item \textbf{Khối lượng tác vụ:} Hàng ngàn lượt đặt giá và cập nhật trạng thái phiên mỗi ngày.
        \item \textbf{Khó khăn thực tế từ người tham gia:}
        \begin{itemize}
            \item \textit{Người mua:} Bị ``bắn tỉa cướp giá'' ở giây cuối, không thể canh màn hình liên tục.
            \item \textit{Người bán:} Bị người thắng bỏ chạy, không trả tiền làm đóng băng hàng hóa.
        \end{itemize}
        \vspace{0.2cm}
        \item \textbf{Nguyên tắc giải quyết nghiệp vụ độc lập với công nghệ:}
        \begin{alertblock}{Cần chuẩn hóa quy chế trước khi phần mềm hóa}
            \begin{itemize}
                \item Không thu cọc trước để tăng độ thanh khoản; áp dụng quy chế nghiêm khắc: \textbf{``3 gậy bùng hàng = Khóa vĩnh viễn tài khoản''}.
                \item Áp dụng quy tắc \textbf{Giá sàn bảo lưu ẩn} để bảo vệ người bán không bị bán lỗ.
            \end{itemize}
        \end{alertblock}
    \end{itemize}
\end{frame}

% ==================== SLIDE 8 ====================
\begin{frame}{Những câu hỏi cốt lõi khi khảo sát sàn đấu giá}
    \begin{columns}[t]
        \begin{column}{0.24\textwidth}
            \begin{block}{\centering \textbf{Làm gì?}}
                \footnotesize Sàn đấu giá tăng dần, tự động nâng giá và giữ tiền bảo chứng an toàn.
            \end{block}
        \end{column}
        \begin{column}{0.24\textwidth}
            \begin{block}{\centering \textbf{Ai, ở đâu?}}
                \footnotesize Người mua \& Người bán giao dịch trực tuyến; Ban quản trị điều hành tập trung.
            \end{block}
        \end{column}
        \begin{column}{0.24\textwidth}
            \begin{block}{\centering \textbf{Dữ liệu nào?}}
                \footnotesize Món hàng, Giá trần bí mật, Lịch sử trả giá công khai, Số dư ví, Đơn hàng.
            \end{block}
        \end{column}
        \begin{column}{0.24\textwidth}
            \begin{block}{\centering \textbf{Thách thức?}}
                \footnotesize Tranh chấp xung đột ở mili-giây cuối, đòi hỏi kỹ thuật khóa dòng dữ liệu.
            \end{block}
        \end{column}
    \end{columns}
\end{frame}

% ==================== SLIDE 9 ====================
\begin{frame}
    \begin{center}
        \vspace{1.5cm}
        {\Large \textbf{\color{primaryNavy}PHẦN 2}}\\[0.5cm]
        {\huge \textbf{\color{primaryNavy}Ví dụ minh họa hiện trạng nghiệp vụ}}\\[0.6cm]
        \textit{\color{secondarySlate}Hai quy trình cốt lõi: Đấu giá ủy quyền tự động \& Chốt đơn Ví bảo chứng}
    \end{center}
\end{frame}

% ==================== SLIDE 10 ====================
\begin{frame}{Sơ đồ hiện trạng: Đấu giá ủy quyền tự động (Proxy Bidding -- Tự động nâng giá)}
    \begin{center}
    \resizebox{0.92\textwidth}{!}{
    \begin{tikzpicture}[
        >=stealth,
        role/.style={rectangle, draw=primaryNavy, fill=primaryNavy, text=white, font=\bfseries\small, minimum width=3.8cm, minimum height=0.65cm, rounded corners=2pt, align=center},
        stepbox/.style={rectangle, draw=primaryNavy!30, fill=primaryNavy!85, text=white, font=\footnotesize, rounded corners=3pt, minimum width=3.6cm, minimum height=0.8cm, align=center},
        alertbox/.style={rectangle, draw=red!40, fill=red!75!primaryNavy, text=white, font=\footnotesize, rounded corners=3pt, minimum width=3.6cm, minimum height=0.8cm, align=center},
        livebox/.style={rectangle, draw=teal!40, fill=teal!70!primaryNavy, text=white, font=\footnotesize, rounded corners=3pt, minimum width=3.6cm, minimum height=0.8cm, align=center}
    ]
        % --- NỀN CÁC LÀN BƠI HIỆN TRẠNG (SWIMLANES) ---
        \fill[black!3]        (-0.1, 0.45) rectangle (3.9, -4.8);
        \fill[primaryNavy!3]   (4.1, 0.45) rectangle (8.1, -4.8);
        \fill[black!3]        (8.3, 0.45) rectangle (12.3, -4.8);

        % --- ĐƯỜNG PHÂN CÁCH LÀN (DASHED DIVIDERS) ---
        \draw[dashed, gray!40, thick] (4.0, 0.45) -- (4.0, -4.8);
        \draw[dashed, gray!40, thick] (8.2, 0.45) -- (8.2, -4.8);

        % --- TIÊU ĐỀ 3 VAI TRÒ CHỊU TRÁCH NHIỆM ---
        \node[role] (r1) at (1.9, 0.0) {Người mua mới (Bidder)};
        \node[role] (r2) at (6.1, 0.0) {Bộ máy xử lý Sàn (System)};
        \node[role] (r3) at (10.3, 0.0) {Người dẫn đầu cũ (Outbid)};

        % --- CÁC BƯỚC DIỄN TIẾN BẬC THANG THEO THỜI GIAN ---
        % Bước 1: Người mua mới đặt giá trần bí mật (Max Bid)
        \node[stepbox] (s1) at (1.9, -1.0) {
            \textbf{1. Đặt giá trần bí mật (Max Bid)}\\[2pt]
            \scriptsize Tối đa: 16.000.000đ $\rightarrow$ Gửi lệnh
        };

        % Bước 2: Sàn khóa phiên & so khớp tự động
        \node[stepbox] (s2) at (6.1, -2.1) {
            \textbf{2. Khóa phiên, so khớp tự động}\\[2pt]
            \scriptsize Tự động nhảy giá lên \textbf{14.200.000đ}
        };

        % Bước 3: Thông báo đẩy cho người cũ bị vượt giá (Outbid Alert)
        \node[alertbox] (s3) at (10.3, -3.2) {
            \textbf{3. Bị vượt giá! (Outbid Alert)}\\[2pt]
            \scriptsize (Đã bị vượt, có thể nâng Max Bid)
        };

        % Bước 4: Phát sóng WebSocket cập nhật toàn sàn (Broadcast)
        \node[livebox] (s4) at (6.1, -4.2) {
            \textbf{4. Phát sóng toàn sàn (Broadcast)}\\[2pt]
            \scriptsize WebSocket đồng bộ giá mới tức thời
        };

        % --- MŨI TÊN KẾT NỐI BẬC THANG CHUẨN ZÍCH ZẮC (-| VÀ --) ---
        % s1 -> s2: Đi ngang sang cột 2 rồi cắm thẳng vào đỉnh s2
        \draw[->, thick, primaryNavy!85] (s1.east) -| (s2.north)
            node[above, pos=0.35, font=\scriptsize\bfseries\color{primaryNavy}] {Gửi giá (Place Bid)};

        % s2 -> s3: Đi ngang sang cột 3 rồi cắm thẳng vào đỉnh s3
        \draw[->, thick, primaryNavy!85] (s2.east) -| (s3.north)
            node[above, pos=0.35, font=\scriptsize\bfseries\color{primaryNavy}] {Báo vượt (Outbid)};

        % s2 -> s4: Đi thẳng xuống s4 trong cùng cột Sàn
        \draw[->, thick, primaryNavy!85] (s2.south) -- (s4.north)
            node[right, pos=0.5, font=\scriptsize\bfseries\color{primaryNavy}] {Đồng bộ (Broadcast)};

    \end{tikzpicture}
    }
    \end{center}
    \vspace{0.05cm}
    \begin{center}
        \textit{\small\color{secondarySlate}Mỗi bước gắn với đúng một vai trò chịu trách nhiệm --- đây chính là ``sơ đồ hiện trạng'' kinh điển trong khảo sát nghiệp vụ.}
    \end{center}
\end{frame}

% ==================== SLIDE 11 ====================
\begin{frame}{Sơ đồ hiện trạng: Chốt phiên, Ví bảo chứng \& Giải ngân}
    \begin{center}
    \resizebox{0.92\textwidth}{!}{
    \begin{tikzpicture}[
        >=stealth,
        role/.style={rectangle, draw=primaryNavy, fill=primaryNavy, text=white, font=\bfseries\small, minimum width=3.8cm, minimum height=0.65cm, rounded corners=2pt, align=center},
        stepbox/.style={rectangle, draw=primaryNavy!30, fill=primaryNavy!85, text=white, font=\footnotesize, rounded corners=3pt, minimum width=3.6cm, minimum height=0.8cm, align=center},
        successbox/.style={rectangle, draw=teal!40, fill=teal!70!primaryNavy, text=white, font=\footnotesize, rounded corners=3pt, minimum width=3.6cm, minimum height=0.8cm, align=center}
    ]
        % --- NỀN CÁC LÀN BƠI HIỆN TRẠNG (SWIMLANES) ---
        \fill[black!3]        (-0.1, 0.45) rectangle (3.9, -5.6);
        \fill[primaryNavy!3]   (4.1, 0.45) rectangle (8.1, -5.6);
        \fill[black!3]        (8.3, 0.45) rectangle (12.3, -5.6);

        % --- ĐƯỜNG PHÂN CÁCH LÀN (DASHED DIVIDERS) ---
        \draw[dashed, gray!40, thick] (4.0, 0.45) -- (4.0, -5.6);
        \draw[dashed, gray!40, thick] (8.2, 0.45) -- (8.2, -5.6);

        % --- TIÊU ĐỀ 3 VAI TRÒ CHỊU TRÁCH NHIỆM ---
        \node[role] (r1) at (1.9, 0.0) {Người mua (Buyer)};
        \node[role] (r2) at (6.1, 0.0) {Hệ thống Sàn (Escrow)};
        \node[role] (r3) at (10.3, 0.0) {Người bán (Seller)};

        % --- CÁC BƯỚC BẬC THANG THEO THỜI GIAN ---
        % Bước 1: Sàn chốt phiên & sinh đơn
        \node[stepbox] (s1) at (6.1, -0.9) {
            \textbf{1. Hết giờ $\rightarrow$ Đạt giá sàn}\\[2pt]
            \scriptsize Tự động lập \textbf{Đơn hàng trúng giá}
        };

        % Bước 2: Người mua thanh toán ví bảo chứng
        \node[stepbox] (s2) at (1.9, -2.0) {
            \textbf{2. Thanh toán đơn hàng}\\[2pt]
            \scriptsize Tiền tạm giữ vào \textbf{Ví bảo chứng}
        };

        % Bước 3: Người bán giao hàng
        \node[stepbox] (s3) at (10.3, -3.1) {
            \textbf{3. Nhận thông báo tiền đã giữ}\\[2pt]
            \scriptsize Đóng gói \& Giao cho vận chuyển
        };

        % Bước 4: Người mua nhận hàng & xác nhận
        \node[stepbox] (s4) at (1.9, -4.2) {
            \textbf{4. Kiểm tra hàng $\rightarrow$ Xác nhận}\\[2pt]
            \scriptsize Bấm ``Đã nhận hàng thành công''
        };

        % Bước 5: Sàn giải ngân & thu phí
        \node[successbox] (s5) at (6.1, -5.1) {
            \textbf{5. Trích thu phí sàn \& Giải ngân}\\[2pt]
            \scriptsize Thu 5\% phí $\rightarrow$ Giải ngân 95\% cho Seller
        };

        % --- MŨI TÊN KẾT NỐI BẬC THANG CHUẨN ZÍCH ZẮC ---
        \draw[->, thick, primaryNavy!85] (s1.west) -| (s2.north)
            node[above, pos=0.35, font=\scriptsize\bfseries\color{primaryNavy}] {Báo nộp tiền};

        \draw[->, thick, primaryNavy!85] (s2.east) -| (s3.north)
            node[above, pos=0.35, font=\scriptsize\bfseries\color{primaryNavy}] {Báo giao hàng};

        \draw[->, thick, primaryNavy!85] (s3.west) -| (s4.north)
            node[above, pos=0.35, font=\scriptsize\bfseries\color{primaryNavy}] {Phát hàng};

        \draw[->, thick, primaryNavy!85] (s4.east) -| (s5.north)
            node[above, pos=0.35, font=\scriptsize\bfseries\color{primaryNavy}] {Xác nhận nhận hàng};

    \end{tikzpicture}
    }
    \end{center}
    \vspace{0.05cm}
    \begin{center}
        \textit{\small\color{secondarySlate}Cơ chế ví bảo chứng thanh toán: Người mua nhận đúng hàng mới giải ngân, loại trừ 100\% rủi ro lừa đảo.}
    \end{center}
\end{frame}

% ==================== SLIDE 12 ====================
\begin{frame}{Bảng mô tả hồ sơ quản lý đấu giá \& đơn hàng}
    \small
    \begin{table}[h]
        \centering
        \begin{tabular}{c l p{7.5cm}}
            \toprule
            \textbf{Mã HS} & \textbf{Tên hồ sơ / Dữ liệu} & \textbf{Vai trò \& Mục đích sử dụng} \\
            \midrule
            \textbf{D1} & Phiếu Món hàng & Lưu thông tin sản phẩm, danh mục, tình trạng, album ảnh. \\
            \textbf{D2} & Phiếu Phiên đấu giá & Thông số: Giá khởi điểm, bước giá, giá sàn ẩn, giá mua ngay. \\
            \textbf{D3} & Lệnh đặt giá trần & Mức giá tối đa bí mật mà người mua cài đặt để hệ thống tự động trả giá thay. \\
            \textbf{D4} & Sổ nhật ký bước giá & Nhật ký lịch sử các bước nhảy giá công khai cho toàn sàn. \\
            \textbf{D5} & Đơn hàng trúng giá & Xác nhận kết quả, địa chỉ nhận hàng, thời hạn thanh toán 72 giờ. \\
            \textbf{D6} & Chứng từ giao dịch ví & Lưu biến động tài chính: nạp tiền, thanh toán, ký quỹ, giải ngân, phí sàn. \\
            \bottomrule
        \end{tabular}
    \end{table}
\end{frame}

% ==================== SLIDE 13 ====================
\begin{frame}{Bảng mô tả công việc là gì?}
    \begin{itemize}
        \item Liệt kê chi tiết các thông tin liên quan đến từng tác vụ trong sơ đồ hiện trạng.
        \item \textbf{Tên công việc:} Đặt tên chính xác, rõ nghĩa (\textit{So khớp giá tự động, Quét chốt phiên, Thanh toán đơn hàng}).
        \item \textbf{Điều kiện kích hoạt:} Hành động người dùng bấm nút hoặc sự kiện hẹn giờ tự động của hệ thống.
        \item \textbf{Tần suất \& thời lượng:} Đo lường tải hệ thống theo giá trị trung bình (giao dịch/phút, thời gian phản hồi $< 100\text{ms}$).
        \item \textbf{Quy tắc chi phối:} Các quy định về giá sàn, bước giá theo bậc, thời hạn 72h thanh toán.
    \end{itemize}
    \vspace{0.2cm}
    \textit{\color{secondarySlate}Bảng mô tả công việc là cầu nối trực tiếp để xây dựng các lớp Dịch vụ nghiệp vụ (Service) trong mã nguồn.}
\end{frame}

% ==================== SLIDE 14 ====================
\begin{frame}{Ví dụ: Bảng mô tả công việc (Hệ thống đấu giá)}
    \small
    \begin{table}[h]
        \centering
        \begin{tabular}{p{3.2cm} p{3cm} p{2cm} p{4.6cm}}
            \toprule
            \textbf{Công việc} & \textbf{Điều kiện kích hoạt} & \textbf{Tần suất} & \textbf{Quy tắc liên quan} \\
            \midrule
            \textbf{So khớp giá đấu tự động} & Có người mua gửi mức giá trần mới & 500 -- 2.000 lượt/ngày & Bước giá theo bậc quy định; Người bán không được tự đấu giá hàng của mình \\
            \midrule
            \textbf{Tự động đóng phiên \& lập đơn} & Đồng hồ điểm hết giờ hoặc có lệnh Mua ngay & Quét mỗi 5 giây & Kiểm tra đạt giá sàn ẩn; Cấp thời hạn thanh toán 72 giờ \\
            \midrule
            \textbf{Giải ngân tiền bán hàng} & Người mua bấm ``Đã nhận hàng'' & 50 -- 200 đơn/ngày & Phí sàn 5\% chuyển về doanh thu; 95\% chuyển về ví Người bán \\
            \bottomrule
        \end{tabular}
    \end{table}
\end{frame}

% ==================== SLIDE 15 ====================
\begin{frame}{THẢO LUẬN 1: Phân tích quy trình Mua đứt ngay}
    \begin{block}{Nội dung thảo luận}
        \textit{Mô tả luồng xử lý khi người mua quyết định sở hữu món hàng ngay lập tức mà không cần chờ phiên đếm ngược hết giờ.}
    \end{block}
    \vspace{0.3cm}
    \textbf{Gợi ý hướng phân tích luồng:}
    \begin{itemize}
        \item \textbf{Bước 1 -- Người mua:} Bấm chọn nút ``Mua ngay'' với mức giá cố định được niêm yết sẵn.
        \item \textbf{Bước 2 -- Hệ thống sàn:} Kiểm tra điều kiện hợp lệ $\rightarrow$ Lập tức khóa phiên $\rightarrow$ Đổi trạng thái sang \textbf{ĐÃ BÁN THÀNH CÔNG} $\rightarrow$ Khởi tạo Đơn hàng.
        \item \textbf{Bước 3 -- Các người tham gia khác:} Nút Mua ngay lập tức biến mất; màn hình cập nhật trạng thái đã bán thành công theo thời gian thực.
    \end{itemize}
\end{frame}

% ==================== SLIDE 16 ====================
\begin{frame}
    \begin{center}
        \vspace{1.5cm}
        {\Large \textbf{\color{primaryNavy}PHẦN 3}}\\[0.5cm]
        {\huge \textbf{\color{primaryNavy}Hiện trạng Tin học}}\\[0.6cm]
        \textit{\color{secondarySlate}Phần cứng, phần mềm và năng lực con người đáp ứng triển khai hệ thống}
    \end{center}
\end{frame}

% ==================== SLIDE 17 ====================
\begin{frame}{Ba mảng khảo sát Tin học}
    \begin{columns}[t]
        \begin{column}{0.32\textwidth}
            \begin{block}{\centering \textbf{Phần cứng}}
                \begin{itemize}
                    \item Máy chủ Web và Cơ sở dữ liệu tiêu chuẩn, đường truyền Internet tốc độ cao.
                    \item Người dùng truy cập qua máy tính cá nhân hoặc điện thoại thông minh.
                \end{itemize}
            \end{block}
        \end{column}
        \begin{column}{0.32\textwidth}
            \begin{block}{\centering \textbf{Phần mềm}}
                \begin{itemize}
                    \item Nền tảng: \textbf{Java 17}, \textbf{Spring Boot 3}.
                    \item Cơ sở dữ liệu: \textbf{MySQL 8.0}.
                    \item Giao diện: Bootstrap 5 kết hợp truyền dữ liệu thời gian thực WebSocket.
                \end{itemize}
            \end{block}
        \end{column}
        \begin{column}{0.32\textwidth}
            \begin{block}{\centering \textbf{Con người}}
                \begin{itemize}
                    \item Đội ngũ phát triển: Nắm vững kiến trúc khối mô-đun, kỹ thuật khóa dòng dữ liệu.
                    \item Khách hàng: Đã quen thuộc với mua sắm online và thanh toán ví điện tử.
                \end{itemize}
            \end{block}
        \end{column}
    \end{columns}
\end{frame}

% ==================== SLIDE 18 ====================
\begin{frame}
    \begin{center}
        \vspace{1.5cm}
        {\Large \textbf{\color{primaryNavy}PHẦN 4}}\\[0.5cm]
        {\huge \textbf{\color{primaryNavy}Xác định và phân loại yêu cầu}}\\[0.6cm]
        \textit{\color{secondarySlate}Bước chuyển từ khảo sát hiện trạng sang đặc tả yêu cầu hệ thống}
    \end{center}
\end{frame}

% ==================== SLIDE 19 ====================
\begin{frame}{Yêu cầu hệ thống là gì?}
    \begin{itemize}
        \item Một yêu cầu diễn tả một \textbf{chức năng} hệ thống phải thực hiện (\textit{ví dụ: Tự động nhảy bước giá}) hoặc một \textbf{đặc điểm/ràng buộc} mà hệ thống phải có (\textit{ví dụ: Ẩn tuyệt đối giá sàn}).
        \item Yêu cầu chính là \textbf{``cầu nối ngôn ngữ''} giữa mong muốn của bài toán kinh doanh và mã nguồn đội phát triển sẽ xây dựng.
    \end{itemize}
    \vspace{0.4cm}
    \begin{center}
        \colorbox{primaryNavy!10}{\textbf{Khảo sát thực tế} $\longrightarrow$ \textbf{Đặc tả yêu cầu phần mềm} $\longrightarrow$ \textbf{Thiết kế Hướng đối tượng}}
    \end{center}
\end{frame}

% ==================== SLIDE 20 ====================
\begin{frame}{Phân loại yêu cầu hệ thống}
    \begin{columns}[t]
        \begin{column}{0.48\textwidth}
            \begin{block}{\textbf{Yêu cầu chức năng}}
                \begin{itemize}
                    \item \textbf{Lưu trữ:} Thông tin hàng hóa, album ảnh, số dư ví, tài khoản.
                    \item \textbf{Tra cứu:} Tìm kiếm món hàng, lọc theo danh mục, theo dõi giá.
                    \item \textbf{Tính toán:} Thuật toán đấu giá tự động, trích thu phí sàn 5\%.
                    \item \textbf{Kết xuất:} Báo cáo bảng điều khiển quản trị, đẩy giá tức thời qua mạng.
                \end{itemize}
            \end{block}
        \end{column}
        \begin{column}{0.48\textwidth}
            \begin{block}{\textbf{Yêu cầu phi chức năng}}
                \begin{itemize}
                    \item \textbf{Bảo mật:} Ẩn danh tính người trả giá, che giấu giá sàn bảo lưu.
                    \item \textbf{Toàn vẹn:} Số dư không bao giờ âm, định dạng số nguyên lớn cho tiền tệ.
                    \item \textbf{Hiệu năng:} Khóa bi quan chống tranh chấp dữ liệu ở giây cuối.
                    \item \textbf{Khả dụng:} Giao diện trực quan, đồng hồ đếm ngược trực tiếp.
                \end{itemize}
            \end{block}
        \end{column}
    \end{columns}
\end{frame}

% ==================== SLIDE 21 ====================
\begin{frame}{Yêu cầu phi chức năng trong Hệ thống Đấu giá}
    \begin{itemize}
        \item \textbf{Yêu cầu bảo mật dữ liệu:}
        \begin{itemize}
            \item Mật khẩu được mã hóa an toàn theo chuẩn quốc tế.
            \item Tuyệt đối giấu kín mức giá sàn bảo lưu và giá trần tối đa trong dữ liệu gửi về giao diện.
        \end{itemize}
        \item \textbf{Yêu cầu an toàn tài chính:}
        \begin{itemize}
            \item Dùng kiểu số nguyên lớn cho tiền VNĐ (triệt tiêu hoàn toàn sai số làm tròn số thập phân).
            \item Cơ chế khóa chống trùng lặp giúp không bao giờ bị trừ tiền hai lần khi bấm đúp chuột.
        \end{itemize}
        \item \textbf{Yêu cầu hiệu năng \& xử lý đồng thời:}
        \begin{itemize}
            \item Áp dụng kỹ thuật khóa dòng dữ liệu trực tiếp trong CSDL để xếp hàng xử lý các lệnh đặt giá dồn dập ở giây cuối cùng.
        \end{itemize}
    \end{itemize}
\end{frame}

% ==================== SLIDE 22 ====================
\begin{frame}{Tài liệu Đặc tả Yêu cầu Hệ thống}
    \begin{itemize}
        \item Là văn bản mô tả đầy đủ các yêu cầu chức năng và \textbf{25 Quy tắc nghiệp vụ} của hệ thống.
        \item Là căn cứ chung để Chuyên viên phân tích, Thiết kế giao diện, Lập trình viên và Kiểm thử viên cùng thống nhất phạm vi.
        \item Đầu vào bắt buộc cho các bước mô hình hóa hướng đối tượng: \textit{Sơ đồ ca sử dụng, Sơ đồ hoạt động, Sơ đồ tuần tự, Sơ đồ lớp}.
        \item \textbf{Mục đích cốt lõi:} Xác định rõ ranh giới trách nhiệm giữa Người mua, Người bán và Hệ thống Sàn.
    \end{itemize}
\end{frame}

% ==================== SLIDE 23 ====================
\begin{frame}{Từ yêu cầu thu thập đến yêu cầu tốt}
    \begin{columns}[t]
        \begin{column}{0.32\textwidth}
            \begin{block}{\centering \textbf{1. Phân loại \& Chuẩn hóa}}
                Mã hóa yêu cầu rõ ràng, có thể kiểm chứng và đo lường độc lập.
            \end{block}
        \end{column}
        \begin{column}{0.32\textwidth}
            \begin{block}{\centering \textbf{2. Thu thập \& Xác nhận}}
                Tham khảo từ các sàn quốc tế; xác nhận kỹ thuật đóng phiên và chế tài phạt bùng hàng.
            \end{block}
        \end{column}
        \begin{column}{0.32\textwidth}
            \begin{block}{\centering \textbf{3. Hoàn thiện \& Ưu tiên}}
                Loại bỏ mâu thuẫn; xây dựng \textbf{41 bài kiểm thử tự động} để kiểm chứng 100\% yêu cầu.
            \end{block}
        \end{column}
    \end{columns}
\end{frame}

% ==================== SLIDE 24 ====================
\begin{frame}{THẢO LUẬN 2: Phân loại yêu cầu cho Hệ thống Đấu giá}
    \begin{block}{Đề bài thảo luận}
        \textit{Liệt kê 3 yêu cầu chức năng và 2 yêu cầu phi chức năng tiêu biểu cho sàn đấu giá.}
    \end{block}
    \vspace{0.2cm}
    \begin{columns}[t]
        \begin{column}{0.5\textwidth}
            \textbf{3 Yêu cầu chức năng:}
            \begin{enumerate}
                \item Cho phép cài đặt giá trần và tự động nhảy bước giá trả thay người dùng.
                \item Cho phép mua đứt ngay món hàng với mức giá niêm yết sẵn.
                \item Cho phép Người bán đăng lại sản phẩm chưa bán được chỉ với 1 thao tác bấm.
            \end{enumerate}
        \end{column}
        \begin{column}{0.48\textwidth}
            \textbf{2 Yêu cầu phi chức năng:}
            \begin{enumerate}
                \item \textbf{Bảo mật:} Ẩn hoàn toàn giá sàn bảo lưu trong toàn bộ dữ liệu trả về giao diện.
                \item \textbf{Thời gian thực:} Đồng bộ bước giá mới tới tất cả người xem trong dưới 1 giây mà không cần tải lại trang.
            \end{enumerate}
        \end{column}
    \end{columns}
\end{frame}

% ==================== SLIDE 25 ====================
\begin{frame}{Tổng kết Phần 1}
    \begin{itemize}
        \item Khảo sát hiện trạng gồm 3 mảng: \textbf{Tổ chức, Nghiệp vụ và Tin học}.
        \item \textbf{Sơ đồ hiện trạng} mô tả từng bước công việc gắn liền với đúng vai trò chịu trách nhiệm (Người mua, Người bán, Sàn).
        \item \textbf{Bảng mô tả công việc \& hồ sơ} là nền tảng trực tiếp để chuẩn hóa 15 bảng cơ sở dữ liệu và các thực thể dữ liệu trong mã nguồn.
        \item Yêu cầu hệ thống được phân định rành mạch thành \textbf{Yêu cầu chức năng} và \textbf{Yêu cầu phi chức năng}.
        \item Toàn bộ 41 yêu cầu nghiệp vụ đã được kiểm chứng tự động qua bộ bài kiểm thử tự động.
    \end{itemize}
    \vspace{0.3cm}
    \begin{center}
        \textit{\color{secondarySlate}Chương tiếp theo: Kỹ thuật mô hình hóa Hướng đối tượng (UML \& Mẫu thiết kế)}
    \end{center}
\end{frame}

% ==================== SLIDE 26 ====================
\begin{frame}
    \begin{center}
        \vspace{1.5cm}
        {\huge \textbf{\color{primaryNavy}CÂU HỎI \& THẢO LUẬN}}\\[0.8cm]
        {\large \textbf{Cảm ơn cô và các bạn đã chú ý lắng nghe!}}\\[0.5cm]
        \textit{\color{secondarySlate}Nhóm 01 sẵn sàng giải đáp thắc mắc về hiện trạng \& yêu cầu hệ thống.}
    \end{center}
\end{frame}

\end{document}
